%%%PAGE 286 rot=0 legibility=good summary=粘贴的倾斜传送带连接体题（A、B 两物块，求加速度/最远距离及传送带做功/最小热量）及手写解答；叠放方块受力、最速降线、竖直圆周两球、变力做功等零散笔记（页面上缘被裁）
%%%BLOCK id=286-01 ch=03 sec=传送带·倾斜传送带连接体 kind=problem alt=06,01 cont=none y=0.00-0.65
% 题干为粘贴的印刷题条（紫色底、略倾斜，上缘及右缘被裁）；手写在题干中对“v=2m/s”“m1=5kg”“m2=5kg 的小物块B”“v0=8m/s”“μ=0.5”画了圈，对“小物块A”“不可伸长”“此时物块A、B的速率相等，且轻绳绷紧”“从传送带下端冲上传送带”“传送带方向的初速度”“整个运动过程中物块B都没有上升到定滑轮处”“此过程中传送带对物块A做的功”等画了下划线
% 题干旁页边另有零散手写草写：“C B D　B D C　B A”“A D　A C D　B C　B D”“(8−v)/20”“−10”，右缘“20 5”“m/s 5”“100 −”
\begin{problem}[\unclear{17}.（13分）]
如图所示，一足够长的倾斜传送带以速度 $v=2\unit{m/s}$ 沿顺时针方向匀速转动，传送带与水平方向的夹角 $\theta=37^\circ$。质量 $m_1=5\unit{kg}$ 的小物块 A 和质量 $m_2=5\unit{kg}$ 的小物块 B 由跨过定滑轮的轻绳连接，A 与定滑轮间的绳子与传送带平行，轻绳足够长且不可伸长。某时刻开始给物块 A 以沿传送带方向的初速度 $v_0=8\unit{m/s}$（此时物块 A、B 的速率相等，且轻绳绷紧），使物块 A 从传送带下端冲上传送带，已知物块 A 与传送带间的动摩擦因数 $\mu=0.5$，不计滑轮的质量与摩擦，整个运动过程中物块 B 都没有上升到定滑轮处。取 $\sin37^\circ=0.6$，$\cos37^\circ=0.8$，重力加速度 $g=10\unit{m/s^2}$，求：

（1）物块 A 刚冲上传送带时加速度的大小及方向；

（2）物块 A 冲上传送带沿传送带向上运动的最远距离及此过程中传送带对物块 A 做的功；

（3）若传送带以不同的速率 $v$（$0<v<v_0$）沿顺时针方向转动，当 $v$ 取多大时，物块 A 沿传送带运动到最\unclear{…}% 粘贴条右缘被裁，句子未完

（题干旁草写）\unclear{C B D　B D C　B A；A D　A C D　B C　B D}\quad $\dfrac{8-v}{20}$\quad \unclear{$-10$}\quad \unclear{20 5；m/s 5；100 $-$}

%FIG p286_a 0.04 0.303 0.39 0.51
\notefig[0.4]{figures/p286_a.png}

\begin{solution}
（1）A：$m_1g\sin\theta+T+\mu m_1g\cos\theta=m_1a_1$

I阶段\quad B：$m_2g-T=m_2a_1$\qquad $a_1=10\unit{m/s^2}$，方向沿传送带向下

（2）A：$m_1g\sin\theta+T-\mu m_1g\cos\theta=m_1a_2$

（匀减速，II阶段）\quad B：$m_2g-T=m_2a_2$\qquad $a_2=6\unit{m/s^2}$

$x_m=x_1+x_2$\qquad $x_1=\dfrac{v_0^2-v^2}{2a_1}$\qquad $x_2=\dfrac{v^2}{2a_2}$\qquad $x_m=\dfrac{10}{3}\unit{m}$

$W_f=-fx_1+fx_2$\qquad $f=\mu m_1g\cos\theta$\qquad $W_f=-\dfrac{160}{3}\unit{J}$。

（3）$x_{1\text{相}}=x_{1\text{物}}-x_{1\text{传}}=\dfrac{v_0^2-v^2}{2a_1}-v\,\dfrac{v_0-\unclear{v}}{a_1}$\qquad $Q_1=fx_{1\text{相}}=v^2-16v+64\ (\unit{J})$% CHECK: 此处分子手写像“v_0−4v”，按上下文疑为 v_0−v

共速后 $x_{2\text{相}}=x_{2\text{传}}-x_{2\text{物}}=v\,\dfrac{v}{a_2}-\dfrac{v^2}{2a_2}=\dfrac{5}{3}v^2$% CHECK: 末项 5/3·v^2 实为 Q_2=f x_{2相} 之值（x_{2相} 本身应为 v^2/(2a_2)），手写未写 Q_2

$Q=Q_1+Q_2=\dfrac{8}{3}v^2-16v+64\ (\unit{J})$\qquad $\therefore Q_{\min}=40\unit{J}$。

$v=3\unit{m/s}$ 时
\end{solution}
\end{problem}
%%%END
%%%BLOCK id=286-02 ch=02 sec=叠放物体·受力分析 kind=note alt=none cont=none y=0.63-0.71
%FIG p286_b 0.09 0.635 0.205 0.71
\notefig[0.25]{figures/p286_b.png}

$F_{2\text{地}}=2.5G$\quad 两边不算
%%%END
%%%BLOCK id=286-03 ch=06 sec=功能关系·系统内力做功 kind=note alt=03 cont=none y=0.68-0.71
对系统，内力对 A、B 做功，绝对值相等
%%%END
%%%BLOCK id=286-04 ch=03 sec=最速降线 kind=note alt=01,06 cont=none y=0.73-0.85
%FIG p286_c 0.085 0.73 0.375 0.835
\notefig[0.3]{figures/p286_c.png}

$t_{NQ}>t_{MQ}$

$V_{y\max}$\unclear{时}，$F_y=0$\quad $F_{\text{合}}=ma_x$\quad $a$ 沿水平\unclear{方向}
%%%END
%%%BLOCK id=286-05 ch=04 sec=竖直面圆周·临界 kind=note alt=06 cont=none y=0.85-0.96
% 图中标注：顶部小球 V1=√(gr)，底部小球 V2=√(5gr)
%FIG p286_d 0.04 0.855 0.215 0.96
\notefig[0.25]{figures/p286_d.png}

不可能同时位于水平直径两端

\unclear{系统 $E_{p\text{低}}>2mgr$}\quad 不会碰撞
%%%END
%%%BLOCK id=286-06 ch=06 sec=功·变力做功 kind=formula alt=none cont=none y=0.94-0.99
有线性变力\quad $W_F=\dfrac{F_1+F_2}{2}x$
%%%END
%%%PAGEEND 286
